non perpendit, sicut pollicetur in actis coram legato sedis apostolicae posse se\footnote{M se posse.} praestare quendam iuristotheologum\footnote{M iuristam theologum, U iuristatheologum.}, plura posse esse\footnote{U esse posse, M om. posse.} capita subordinata. In quo mystici capitis ratio aut etiam symbolici\footnote{M et symbolici, U aut symbolici.} a ratione capitis naturalis deficit. Quare praeter caput Christum et aliud oportere quaeri caput in ecclesia statim probabitur. Neque ei suffragatur, quod apostolum inducit I. ad Cor. 3 [v 5; 1 v. 13]: Divisus est Christus etc.\footnote{fehlt in U.} Nam etsi Paulus ibi Petri meminerit, non tamen falsum dixit beatus Hieronymus libro 1. contra Jovinianum, mihi columna 18: Unus eligitur (loquitur de Petro) ut capite constituto schismatis\footnote{M ut schismatis.} tollatur occasio. Clare appellat Petrum caput in ecclesia constitutum\footnote{U esse constitutum.}. Sed haec missa faciamus\footnote{M U facimus.}, solutiones ad nostra inducta repulsuri\footnote{M refusuri.}.

Primo cum respondit\footnote{M U respondet.} ad illud Johannis 5: Neque enim potest filius\footnote{M U filius potest.} facere quicquam, nisi quae viderit patrem facientem, dixit secundum sanctos patres exprimi hic\footnote{fehlt in M.} aequalitatem patris et filii. Verum legat attentius quaeso reverendus pater beatum et inadulabilem\footnote{M inadulabilem, Mc adulabilem.} patrem Bernhardum libro 3. de consideratione ad Eugenium, ubi de forma ecclesiae loquens et eam esse de iure divino probans, nostram formavit rationem\footnote{M formavit probationem, U rationem formavit.}, mihi columna 7.\footnote{mihi col. VII fehlt in M.} Nec vilem reputes\footnote{U reputans, Uv reputes.} formam hanc, quia in terra est: exemplar habet in coelo. Neque enim filius potest facere quicquam, nisi quae viderit patrem facientem. Praesertim cum ei sub Moysi nomine dictum sit: Vide, omnia\footnote{M ne omnia.} facias secundum exemplar, quod tibi in monte monstratum est. Viderat hoc, qui dicebat\footnote{U dicebar.}: Vidi civitatem sanctam\footnote{U add. etc.}. Ego enim propter similitudinem dictum reor: quod sicut illic Seraphim et Cherubim\footnote{M cherubim et seraphim.} et 'ceteri' quique\footnote{U ceteri qui; M ceteri, constitutique, gestrichen constituti.} usque ad angelos et archangelos ordinantur sub uno capite deo, ita hic quoque sub uno summo pontifice primates vel patriarchae, episcopi, presbyteri\footnote{M Archipresbyteri statt episcopi, presbyteri; U Archiepiscopi, episcopi, presbyteri.} vel abbates et reliqui in hunc modum. Tunc\footnote{fehlt in M.} 'subdit Bernhardus': Non est parvi pendendum, quod et deum habet\footnote{U habeat.} auctorem, et de caelo ducit originem. Quis iam non intelligit\footnote{U intelligat.}, hanc ecclesiasticam hierarchiam secundum Bernhardum a Christo institutam? Et velut Christus\footnote{M U deus.} caput est in caelo, ita summus pontifex caput est\footnote{M U sit.} in militante ecclesia, neutiquam\footnote{M nec dicam.} Christum excludendo, cuius se fatetur vicarium. Sed ad ratiunculam plebeiam\footnote{U pleiam, Uv plebeiam.}, quam inducit\footnote{U induxit.}, ecclesiam manere acephalam\footnote{M ad Scepolam.}, [261] papa mortuo, nisi ridicule diceremus Christum cedere papae vivo et succedere papae mortuo, 'ridicula prorsus ratiuncula', quae\footnote{M et quae.} in re tam seria, inter tot praestantes viros, vix est digna renarrari\footnote{M U enarrari.}, cum ab initio dixerim, caput illud symbolicum a veri capitis 'naturalis'\footnote{M naturalis 2 mal.} ratione in aliquibus convenientiis deficere. Neque Christus, cuius regnum manet in aeternum et sempiternum est eius\footnote{U et cuius sempiternum est.} sacerdotium, cedit papae vel succedit, cum ei sit data omnis potestas in caelo et in terra (Matthaei ultimo) et defuncto papa modo coetus cardinalium, sicut defuncto episcopo capitulum, iura illa teneant\footnote{M U teneat, P Uv teneant.}, usquequo novus pontifex deligatur\footnote{M delegatur.}.

Quod vero secundo loco reverendus pater dicit, Cyprianum loqui de quolibet episcopo, non de Romano pontifice, plurimum miror, cum intelligentia dictorum ex causis dicendi\footnote{M dicendis.} sit accipienda et Cyprianus in locis per me\footnote{M a me.} adductis eos increpet, qui a Cornelio defecerunt, qui certe Romanus fuit pontifex. Ob id sciat reverendus pater me non nudis\footnote{M nullis, Mc nudis.} verbis satiari\footnote{M satiare, Mc satiari.}, ut sophistas pascere solemus. Quae ex Cypriano pro se affert, opponendo arbitror cumulatius efficiet. Nam quod Cyprianus Cornelium fratrem appellat et apostolos fratres fuisse, nemo nescit. Attamen Petrus sicut et eius successor Cornelius caput fuit apostolorum, apex et vertex iuxta beati Dionysii ca. 3. vel. 7. de divinis nominibus sententiam\footnote{fehlt in M (Lücke).}. De electione et a\footnote{fehlt in U (ist in Uv zugefügt).} concilio Niceno et a Cypriano conscripta nec promovent nec impediunt negotium. Et multo minus Augustinus lib. 2. cap. 2. de baptismo parvulorum debebat adduci post Cyprianum. Nam Augustinus arrogantiam et temeritatem eorum, qui per ambitionem et superbiam ad ecclesiasticas praelaturas se ingerunt, reprehendit: quod non debeant se constituere nec alios cogere ad se constituendum, cum quisque praelatus expectare debeat, ut vocetur, sicut Aaron.

Tertio aliam Cypriani sententiam, ex tertia epistola\footnote{P M U ex secunda epistola.} ad Cornelium, diluturus\footnote{M dilaturus.} respondit unitatem quidem sacerdotalem in occidentali ecclesia a Romana exortam, non in orientali. Obticuit hic reverendus pater, quod Cyprianus praecedenter Romanam ecclesiam appellavit cathedram Petri\footnote{M U petri cathedram.} atque ecclesiam principalem. Sed quid sua solutio possit, medullam verborum inspicienti est manifestarium. Nam reverendus pater more grammatico ortum unitatis sacerdotalis intellexit\footnote{M U intelligit.} quo ad auspicia\footnote{U auspicium.} et initium\footnote{M initia.}, cum profecto Cyprianus voluerit explicare ortum commissionis,
